% 01-introduction.tex — Introduction : pourquoi (compréhension des étudiants
% et arrivée de l'IA), pour qui (étudiants de BAC1, enseignants), comment
% (assistance IA encadrée), objectifs et questions, périmètre, plan.

\chapter{Introduction}\label{chap:introduction}

\section{Pourquoi ce travail}

En première année d'université, la mécanique est souvent la première
rencontre des étudiants avec une physique formalisée : des lois écrites sous
forme vectorielle, appliquées à des systèmes qui évoluent dans le temps.
Beaucoup apprennent à appliquer les formules sans se représenter ce qu'elles
décrivent. Un étudiant peut calculer correctement l'accélération d'un objet
en mouvement circulaire et penser malgré tout qu'une \enquote{force
centrifuge} le pousse vers l'extérieur. Ces idées issues du sens commun sont
connues depuis longtemps et résistent à un enseignement classique
\cite{halloun1985common}. Les cours où l'étudiant est actif et confronte ses
idées à ce qui se passe réellement donnent de bien meilleurs résultats
\cite{hake1998} (chapitre~\ref{chap:mecanique}).

Une partie du problème tient aux supports. Un schéma au tableau ou une
figure du syllabus est immobile : il montre le système à un instant, pas la
façon dont la vitesse, l'accélération et les forces évoluent pendant le
mouvement. Une animation interactive comble ce manque. L'étudiant voit
l'objet bouger, les vecteurs changer de direction et de longueur, et peut
modifier un paramètre pour observer son effet. Des simulations conçues de
cette façon ont montré leur intérêt pour l'apprentissage
\cite{wieman2008phet, rutten2012} (chapitre~\ref{chap:simulations}).

Produire de telles animations reste pourtant coûteux. Il faut maîtriser la
physique, pour que le modèle soit juste, et la programmation d'interfaces,
pour qu'il soit lisible et manipulable. C'est ce coût que l'arrivée de
l'intelligence artificielle générative remet en question. Depuis 2021, des
modèles de langage entraînés sur du code écrivent une fonction à partir de
sa description \cite{chen2021codex}, et les assistants actuels lisent un
projet entier, le modifient, lancent les tests et corrigent leurs erreurs
(chapitre~\ref{chap:ia}). Une animation qui demandait des semaines de
développement devient envisageable en quelques jours. Ce mémoire part de la
rencontre de ces deux constats : les étudiants ont besoin de \emph{voir} la
physique, et l'IA pourrait rendre possible de la leur montrer à moindre coût
\cite{sujet_memoire}.

\section{Pour qui}

Le public visé en premier lieu est celui des étudiants de première année du
cours de Physique~1 (LEPL1201) de l'École polytechnique de Louvain. Les
animations reprennent le formulaire et les notations de ce cours
\cite{formulaire_epl1201}, pour que l'étudiant y retrouve exactement les
symboles qu'il rencontre au cours et à l'examen. Le cours de physique de la
Faculté des sciences de la motricité (LFSM1105), qui partage la même
mécanique, peut utiliser les mêmes pages.

L'objectif pour l'étudiant est de faciliter la compréhension en
\emph{montrant} ce que les formules calculent. Chaque animation affiche le
phénomène, toutes les forces qui agissent avec leur direction et leur
grandeur, et les formules du cours écrites trois fois : sous forme
symbolique, avec les valeurs numériques du moment remplacées, puis avec le
résultat (figure~\ref{fig:intro-page}). L'étudiant peut ainsi suivre le lien
entre le calcul qu'il ferait sur papier et ce qu'il voit à l'écran.

\begin{figure}[ht]
  \centering
  \includegraphics[width=\linewidth]{looping-ecran-complet.png}
  \caption{Une page du site (looping avec wagon). À gauche, la scène avec les
    forces ; au centre, le diagramme du corps libre et les énergies ; à
    droite, les paramètres et un graphe en fonction du temps ; en bas, les
    formules avec leurs valeurs substituées. Toutes les vues suivent le même
    instant de la simulation.}
  \label{fig:intro-page}
\end{figure}

Le travail s'adresse aussi aux enseignants. Ce sont eux qui utilisent les
animations au cours ou les proposent en travail autonome, et ce sont eux qui
devraient pouvoir en créer de nouvelles pour leurs propres exercices. Le
site du cours propose déjà deux animations \cite{zoulab_epl1201} ; la
question est de savoir comment en produire davantage, sans y consacrer un
temps démesuré et sans sacrifier la justesse.

\section{Comment}

Nous avons développé un site de simulations de mécanique en nous faisant
assister par l'IA, principalement Claude Code (Anthropic) et ChatGPT Codex
(OpenAI). L'enjeu n'est pas de \enquote{faire écrire du code} par l'IA.
Laissé sans cadre, un assistant produit un code qui a l'air juste mais peut
contenir une formule fausse, une hypothèse cachée ou une structure
impossible à faire évoluer ; une part importante du code généré contient par
exemple des failles que l'utilisateur ne voit pas \cite{pearce2022asleep}. En
physique, une animation fausse est d'autant plus dangereuse qu'elle peut
avoir l'air juste.

Nous avons donc encadré l'IA pour obtenir trois propriétés à la fois :
\begin{itemize}
  \item \textbf{une physique correcte} : les formules exactes du cours sont
        utilisées quand elles existent, les hypothèses du modèle sont écrites
        sur la page, et chaque calcul est vérifié par des tests et par un
        agent séparé qui redérive les formules et les recoupe numériquement ;
  \item \textbf{un code maintenable} : tous les exercices suivent la même
        structure (un fichier de calcul, un fichier d'affichage, une page),
        les calculs physiques ne sont jamais partagés entre exercices pour
        rester faciles à vérifier, et le site n'utilise aucune bibliothèque
        externe ;
  \item \textbf{une création facilitée et en partie automatisée} : les
        règles du projet sont écrites une seule fois dans un fichier lu par
        l'assistant à chaque session, et des procédures réutilisables
        (\emph{skills}) et des agents spécialisés prennent en charge les
        tâches récurrentes : créer un exercice, écrire les tests, vérifier
        la physique, contrôler qu'un exercice est terminé.
\end{itemize}
Les choix pédagogiques, la validation des résultats et la responsabilité de
leur exactitude restent les nôtres. L'IA accélère le travail, le cadre en
garantit la qualité, et ce mémoire évalue les deux. Le site est publié en
ligne \cite{site_physicsviz}.

\section{Objectifs et questions de recherche}

Le sujet fixe deux objectifs \cite{sujet_memoire}. Le premier est de
développer un ensemble d'animations JavaScript interactives pour la physique
de première année, physiquement justes et utiles en classe. Le second est
d'évaluer de façon critique ce que les outils d'IA apportent au
programmeur-concepteur, pour écrire le code comme pour organiser
l'application. Nous en tirons trois questions :
\begin{enumerate}
  \item \textbf{Conception.} Quelles caractéristiques d'une animation aident
        un étudiant de première année à relier le phénomène, les forces et
        les formules du cours ?
  \item \textbf{Développement.} Des enseignants qui programment peuvent-ils
        produire de telles animations avec un assistant d'IA, avec une
        physique vérifiée et un code maintenable ? Que gagnent-ils en temps,
        et à quelle étape ?
  \item \textbf{Limites.} Où l'assistant se trompe-t-il, et quelles
        vérifications faut-il pour qu'une erreur n'atteigne pas les
        étudiants ?
\end{enumerate}
L'état de l'art répond en grande partie à la première question et prépare les
deux autres, qui sont traitées dans les parties Réalisation et Résultats.

\section{Périmètre}

Le site compte aujourd'hui cinq exercices de mécanique : le tir parabolique,
le pendule, le mouvement circulaire uniforme et deux versions du looping.
Quatre d'entre eux proposent un mode jeu expérimental, dans lequel
l'étudiant doit atteindre un objectif en réglant les conditions initiales ;
dans les deux versions du looping, il s'agit de faire tourner un seau d'eau
sans le renverser. Le sujet couvre aussi l'électricité élémentaire, que le
formulaire du cours aborde (force de Coulomb, théorème de Gauss, potentiel
électrique) \cite{formulaire_epl1201}. [À COMPLÉTER : exercices d'électricité
ajoutés d'ici le dépôt, ou électricité hors périmètre, et pourquoi.]

\section{Plan du mémoire}

Le mémoire est organisé en trois parties.

La \textbf{partie~I}, \emph{État de l'art}, rassemble ce que l'on sait déjà.
Le chapitre~\ref{chap:mecanique} résume les difficultés des étudiants en
mécanique. Le chapitre~\ref{chap:simulations} présente ce que la recherche
montre de l'apprentissage avec des simulations, ainsi que les outils
existants. Le chapitre~\ref{chap:ludification} traite des éléments de jeu
dans l'apprentissage. Le chapitre~\ref{chap:numerique} passe en revue les
méthodes de calcul et d'affichage utilisables dans un navigateur. Le
chapitre~\ref{chap:ia} fait le point sur les assistants de programmation et
se termine par les exigences que nous en tirons pour le site.

La \textbf{partie~II}, \emph{Réalisation}, décrit ce que nous avons
construit : nos outils et nos choix techniques
(chapitre~\ref{chap:cahier}), l'architecture du site
(chapitre~\ref{chap:architecture}), les exercices
(chapitre~\ref{chap:exercices}), le mode jeu du seau d'eau
(chapitre~\ref{chap:jeu}), notre façon d'utiliser l'IA avec ses règles,
skills et agents (chapitre~\ref{chap:methodologie}) et la validation de la
physique et du code (chapitre~\ref{chap:validation}).

La \textbf{partie~III}, \emph{Résultats}, présente la façon dont nous
évaluons le travail (chapitre~\ref{chap:protocole}), les résultats mesurés,
dont la comparaison entre la théorie et la simulation et l'apport réel de
l'IA (chapitre~\ref{chap:quantitatif}), les retours des utilisateurs
(chapitre~\ref{chap:qualitatif}) et leur discussion
(chapitre~\ref{chap:discussion}). La conclusion
(chapitre~\ref{chap:conclusion}) répond aux trois questions et ouvre sur la
suite.
